\section{Introduction}~\label{intro}

Robots operating in real-world environments often have incomplete or unreliable information about the current task state due to perception errors~\cite{kaipa2016addressing}, execution failures~\cite{kaelbling1998planning}, and unexpected environmental changes~\cite{kurniawati2016online}. If such uncertainty is not resolved, the robot may make incorrect decisions and fail the task. External information from an expert can help resolve this uncertainty~\cite{dragan2013legibility, goodrich2008human}, but frequent queries can incur substantial costs in time and workload~\cite{chen2014human}. The robot must therefore determine \textit{when} external information is necessary and \textit{what} information should be requested.

Recent uncertainty-aware systems selectively request external feedback during task execution. Existing approaches commonly trigger queries based on prediction confidence, including conformal prediction~\cite{ren2023robots, liang2024introspective, mullen2025lbap}, or on the expected effect of a query on task reward~\cite{burks2021collaborative, anand2025adaptive}. While these strategies reduce unnecessary interaction, they do not explicitly \cp{require confidence in} the current \cp{symbolic} state \cp{to reach a specified level}. As a result, they may fail to identify when additional state information is needed or what information would be most useful to request. In contrast, we determine when to query based on state uncertainty and what to query based on its expected reduction.

In this paper, we present a method that \cp{constructs} a weighted belief over \cp{reachable successor} states \cp{after each action and observation}. We formulate decision making as a partially observable Markov decision process (POMDP)~\cite{kaelbling1998planning} and represent task knowledge using a STRIPS-style symbolic model~\cite{davis1993knowledge,fikes1971strips}. The belief is updated after each action and observation, and its uncertainty is used to determine when external information is needed. If a query is triggered, the robot selects the state fact expected to reduce uncertainty the most and asks an external expert to verify it. The response is incorporated into the belief, which is then used for subsequent planning and execution.


We evaluate the method in tomato harvesting and waste sorting, which exhibit different forms of state uncertainty. Using oracle feedback, we analyze the effects of query timing and query selection and compare against conformal-prediction-based and reward-based strategies. We then evaluate non-oracle feedback from a VLM and human participants and compare against KnowNo~\cite{ren2023robots}. \cp{The planned} interviews \cp{examine how} participants \cp{interpret robot questions and choose responses}.


The main contributions of this work are as follows\cp{.}

\begin{itemize}

\item We determine when to query based on uncertainty in the current task state and show that query timing has a substantial effect on task success.

\item We prioritize state information that is expected to reduce uncertainty and show that query selection affects the number of required queries.

\cp{\item We show higher task success with fewer queries than Query-Action in oracle simulation and higher success with additional queries than KnowNo during physical execution with VLM feedback.}

\end{itemize}
